\section{Geometric propagation and a calibrated path}\label{sec:geo}

\subsection{Ellipse capacity and the propagation exponent}

We parametrize lines by \(P=(U,V)\in\R^2\times\R^2\), and write
\[
X_v=V+vU.
\]
Throughout this section, two tests using the same ellipse may have
independently chosen centers.  Thus \(U\in E,\ V\in E\) abbreviates
\(U\in u+E,\ V\in v+E\), with arbitrary \(u,v\in\R^2\); it does not
require the two centers to coincide.

Fix \(0<\kappa<1/10\).  For a centered ellipse \(E\) with radii
\(L\ge w>0\), define
\[
W(E)=L^{1+\kappa}w^{1-\kappa}.
\]
We call \(W(E)\) its capacity weight.  Rectangles with the same
semiaxes may replace ellipses in upper bounds, at the cost of an
absolute constant.  The same is true of dilations by bounded factors.
In particular, a constant enlargement of a tested ellipse is covered
by boundedly many tests at the original widths.

Let \(M=2^n\), with \(n\to\infty\).  A time test at depth \(y\) uses
the dyadic intervals of length \(2^{-\lfloor yn\rfloor}\).  For
\(T\in[0,1)\), let \(T_y=[T]_y\) denote the left endpoint of its
interval at this depth.  Dyadic roundings change the scales below
by bounded factors only.

\Needspace{4\baselineskip}
\begin{definition}\label{geo:large-sublaw}
A \emph{sublaw} of a probability measure \(\mu\) is a finite measure
\(\nu\le\mu\).  Along a sequence with base \(M\to\infty\), a sublaw
has \emph{large mass} if
\[
\nu(\text{whole space})=M^{-o(1)}.
\]
Normalizing such a sublaw changes every upper mass bound by at most
a factor \(M^{o(1)}\).  We use either its unnormalized masses or its
normalized probabilities, recording this factor when necessary.
\end{definition}

All probability spaces may be taken to be standard Borel.  Auxiliary
labels are retained under restriction and pushforward; only the
displayed spatial and time tests enter the mass bounds.  At each
finite stage, compact coordinate windows are discretized by finite
dyadic grids.  Orientations at widths \(L,w\) are discretized at
spacing comparable to \(w/L\).  A translated ellipse of arbitrary
orientation is covered by boundedly many constant enlargements of
grid cells at these matched scales.  Conversely, each grid cell is
covered by boundedly many continuous tests.  All assignments of
heavy cells can therefore be made from finite lists.

Fix \(s\in(0,1]\), and initially allow a positive error \(\epsilon\)
in the time distribution.  There is a base label determining
\(P=(U,V)\), such that, conditionally on this label, almost surely,
\begin{equation}\label{eq:source-1}
\begin{split}
\max_{I\in\mathcal D_j}
 \Pr(T\in I\mid\text{label})
 &\le M^\epsilon 2^{-sj},\\
\#\{I\in\mathcal D_j:
 \Pr(T\in I\mid\text{label})>0\}
 &\le M^\epsilon 2^{sj},
\qquad 0\le j\le n,
\end{split}
\end{equation}
where \(\mathcal D_j\) is the dyadic partition at depth \(j\).
Testing a sufficiently fine fixed grid of ratios \(j/n\) is
equivalent to testing every depth, with a corresponding arbitrarily
small increase of \(\epsilon\): interpolate the upper weight bound
from the preceding grid depth and the support bound from the
following one.

Suppose that \(P\) remains in a bounded set, uniformly along the
sequence, and that for a fixed finite \(D\ge1\),
\begin{equation}\label{eq:source-2}
\Pr(U\in E,\ V\in E)\le K W(E),
\qquad M^{-D}\le w\le L\le1.
\end{equation}
The constant \(K\) may depend on the instance.  For an exponent \(q\),
the propagation assertion is the existence of a large sublaw on
which, simultaneously for all indicated times, centers, and ellipses,
\begin{equation}\label{eq:source-3}
\Pr(T_1=v,\ U\in E,\ X_v\in M^{-1}E)
\le M^{-q+o(1)}K W(E),
\qquad M^{-(D-1)}\le w\le L\le1.
\end{equation}
Figure~\ref{fig:capacity-tests} depicts the input and terminal
tests.  The shared ellipse orientation imposes no independence
condition on the law of the two spatial variables.
\begin{figure}[tb]
\centering
\begin{tikzpicture}[x=1cm,y=1cm,>=Stealth,
  axes/.style={gray!65,thin},
  capacityellipse/.style={draw=linkblue,fill=linkblue!8,thick}]
  \foreach \x/\name in {0/U,4.2/V,8.4/X_v} {
    \begin{scope}[shift={(\x,0)}]
      \draw[axes,->] (-1.25,0)--(1.35,0);
      \draw[axes,->] (0,-.85)--(0,1.15);
      \node at (0,-1.15) {\(\name\)-plane};
    \end{scope}
  }
  \begin{scope}[shift={(.12,.20)},rotate=25]
    \draw[capacityellipse] (0,0) ellipse (1.00 and .37);
    \draw[linkblue,dashed] (-1,0)--(1,0);
  \end{scope}
  \begin{scope}[shift={(4.05,.26)},rotate=25]
    \draw[capacityellipse] (0,0) ellipse (1.00 and .37);
    \draw[linkblue,dashed] (-1,0)--(1,0);
  \end{scope}
  \begin{scope}[shift={(8.57,.21)},rotate=25]
    \draw[capacityellipse] (0,0) ellipse (.50 and .185);
    \draw[linkblue,dashed] (-.5,0)--(.5,0);
  \end{scope}
  \node at (.05,1.4) {\(E+u_0\)};
  \node at (4.15,1.4) {\(E+v_0\)};
  \node at (8.45,1.4) {\(M^{-1}E+x_0\)};
  \draw[decorate,decoration={brace,mirror,amplitude=4pt}]
    (-1.1,-1.53)--(5.3,-1.53);
  \node at (2.1,-1.98) {initial test};
  \draw[->,thick] (5.8,.3)--(6.8,.3)
    node[midway,above] {\(V+vU\)};
\end{tikzpicture}
\caption{An elliptic capacity test uses the same shape and orientation
in the two spatial coordinates, with independently chosen centers.
The terminal test keeps the \(U\)-ellipse and contracts the
position ellipse by \(M^{-1}\). The picture depicts testing sets,
not a product assumption on the law; the displayed contraction is
schematic.}
\label{fig:capacity-tests}
\end{figure}


Let \(q(\epsilon,D)\) be the supremum of the exponents for which this
assertion holds universally along sequences satisfying
\eqref{eq:source-1} and \eqref{eq:source-2}.  Define
\[
q_*=\lim_{\epsilon\downarrow0}\inf_{1\le D<\infty}q(\epsilon,D).
\]
The limit exists by monotonicity in the admissible time error.
Constants in the coordinate bounds, or bounded changes in the
coordinate ranges, are immaterial by covering and isotropic
normalization.

\Needspace{3\baselineskip}
\begin{theorem}[Geometric propagation]\label{geo:main}
For \(0<s\le1\) and \(0<\kappa<1/10\),
\begin{equation}\label{eq:source-4}
q_*\ge1+2s-10\kappa.
\end{equation}
\end{theorem}

The estimate concerns arbitrary joint laws of \(U,V,T\), with the
base label determining \(U,V\).  Both bounds in
\eqref{eq:source-1} are conditional on that label and hold at every
dyadic depth.  The capacity in \eqref{eq:source-2} is one joint test,
with independent centers and a shared ellipse shape.  The conclusion
allows a dominated sublaw; it need not hold for the original law.
For example, if \(U\) is uniform on a fixed disk, \(V=0\), and \(T\)
is independently uniform, a unit test in the first time bin has mass
of order \(M^{-1}\).  Restricting \(T\) to \([1/2,1)\) removes this
focusing event and gives the corresponding \(M^{-3}W(E)\) bound.

\paragraph{Route through the geometric proof.}
We argue by failure of Theorem~\ref{geo:main}.  The following four stages
organize the argument.
\begin{enumerate}[label=\textup{(\roman*)}]
\item In this section, heavy-box counts and angular exclusion reduce
failure to a calibrated path of ellipse widths and reference axes.
The initial complexity \(D\) may increase as the time error decreases;
the bounded-complexity reduction precedes the local limit.
\item Section~\ref{sec:entropy} encodes that obstruction by finite
sheared states.  Regularization supplies hereditary support and weight
bounds before entropy profiles and their tangents are formed.
\item Section~\ref{sec:projection} derives scalar and joint-coordinate
rate inequalities from the planar incidence theorem, including the
case of correlated time and axis parameters.
\item Section~\ref{sec:aspect} excludes every possible aspect velocity
by a backward comparison whose gain exceeds the calibrated rate.
\end{enumerate}
The heavy ellipses and their calibrated path belong to this
counterfactual failure construction.  The universal conclusion is the
retained-law estimate \eqref{eq:source-3} at every strict exponent
\(q<q_*\).

\subsection{Sequence conventions and conditional time profiles}

We first make the order of limits precise.  A statement valid for
every sequence, with subpower retention and subpower losses, is
uniform with arbitrarily small fixed power losses once the base
is sufficiently large and the other parameters are fixed.
Otherwise, selecting a violating instance at each of arbitrarily
large bases would give a violating sequence.  Conversely, the
fixed-power statements yield subpower losses by a diagonal choice.
This applies simultaneously to any fixed finite collection of
restrictions or estimates.  Passing to subsequences is permitted
throughout.

For clarity, the recurring limit choices are summarized here.  Each
row fixes a finite problem before the underlying base is increased;
the local proofs below give the quantitative choices.
\begin{center}
\small
\begin{tabular}{@{}p{.23\linewidth}p{.72\linewidth}@{}}
\toprule
Stage & Order of choices \\
\midrule
Failure sequence & Fix a positive failure margin, finite complexity,
mesh, and list of restrictions; then increase the base. \\
Calibration & Fix the first-order scale and stretched window, then
choose finite-stage errors negligible in \(\delta\log M\), with
\(\delta\log M\to\infty\). \\
Tangent realization & Fix the tangent radius and finite physical
cutoffs; make previous errors negligible in the new logarithmic unit
before shrinking that radius. \\
Projection tests & Fix a positive output gap, then smaller input
nonconcentration and density losses; finally let the block base tend
to infinity. \\
Backward comparison & Fix the tangent problem and any large velocity;
choose clipping and endpoint errors afterward. \\
\bottomrule
\end{tabular}
\end{center}
In particular, for each \(q<q_*\), the definition of \(q_*\) gives
one sufficiently small time-error threshold valid for every finite
\(D\).  The sufficiently large base threshold may depend on that
fixed \(D\).  We will justify the sequence-to-uniform use of this
statement below and retain this distinction in the packet application.

The input bound immediately gives \(q_*\ge0\).  Indeed, a terminal
test in \eqref{eq:source-3} implies
\(V=X_v-vU\in c+O(E)\), since \(0\le v<1\), and hence its mass is
at most a constant multiple of \(KW(E)\).
Suppose for contradiction that
\[
q_*<1+2s-10\kappa.
\]
All exponents under consideration are then bounded; in particular,
they may be taken less than \(4\).

Choose \(\epsilon_i\downarrow0\), finite complexities \(D_i\), and
exponents tending to \(q_*\) from above, for which the universal
assertion fails.  For each fixed \(i\), the complexity \(D_i\) is
fixed and \(n\to\infty\).  We distinguish errors \(o_i(1)\), which
tend to zero as \(i\to\infty\), from \(o_n(1)\), which tend to zero
along the sequence with \(i\) fixed.

We may choose these sequences with a fixed-power margin at each
fixed \(i\).  More precisely, after leaving room above the relevant
supremum, there are positive tolerances tending to zero with \(i\)
such that even a slightly weaker tested exponent cannot be obtained
while retaining at least the corresponding fixed negative power
of \(M\).  To see this, negate the uniform fixed-power formulation
of the sequence assertion.  If every fixed positive tolerance
worked at every sufficiently large base, a diagonal would give
the sequence assertion itself.  Shrinking a tolerance preserves
failure, so the tolerances can be chosen as small as needed.
Consequently, restrictions retaining \(M^{-o_n(1)}\) mass and changes
of \(M^{o_n(1)}\) in the bounds preserve the chosen failure.  This
also explains why a subsequence on which vanishing-loss conclusions
hold cannot remain inside the selected failing sequence.

\begin{lemma}[Reuse of the time profile]\label{geo:time-conditioning}
Suppose a reference law satisfies \eqref{eq:source-1} with error
\(\epsilon>0\).  Fix a finite collection of prefix depths.  After
any fixed finite succession of restrictions retaining
\(M^{-o_n(1)}\) mass, one may remove a negligible part of the
retained measure so that the following statements hold.

Conditionally on each surviving base label, the retained time law
satisfies \eqref{eq:source-1} with error \(O(\epsilon)\).
Conditionally on a surviving base label and a prefix at depth
\(j_0\), the descendant law has maximum weights and support counts
of exponent \(s\), with error \(O(\epsilon)\) in the original
\(\log M\) units.  Further conditioning on parent cells determined
by the base label and prefix does not change this conclusion.
\end{lemma}

\begin{proof}
All comparisons are made directly with the original reference law.
For a fixed base label, let \(p_I\) be the probability of a prefix
\(I\in\mathcal D_{j_0}\), and let \(N_I(j)\) be its number of
supported descendants at depth \(j\ge j_0\).  The hypotheses give
\[
\max_Ip_I\le M^\epsilon2^{-sj_0},
\qquad
\sum_I N_I(j)\le M^\epsilon2^{sj}.
\]
Thus, when the prefix is sampled according to its original law,
\[
\sum_I p_I N_I(j)
\le M^{2\epsilon}2^{s(j-j_0)}.
\]
For a sufficiently large fixed constant \(C\), Markov's inequality
allows the removal of prefixes for which
\[
N_I(j)>M^{C\epsilon}2^{s(j-j_0)},
\]
at original mass at most \(M^{-(C-2)\epsilon}\).
Summing over \(j\le n\) remains power-small, since \(n=M^{o_n(1)}\).

The prefixes for which
\[
p_I<2^{-sj_0}M^{-C\epsilon}
\]
have total original probability at most
\(M^{-(C-1)\epsilon}\), by the prefix support bound.  First remove
these prefixes and the support-count exceptions at every prescribed
depth.  Then, from the deepest prescribed depth to the shallowest,
discard each label-prefix fiber whose current retained fraction,
relative to its original mass, is less than \(M^{-C\epsilon}\).
Finally apply the same retained-fraction test to the base-label
fibers.  Each retained-fraction stage removes at most
\(M^{-C\epsilon}\) times the original total mass.  A later stage
deletes whole fibers at every
previously checked depth, so the retained fraction in each surviving
fiber keeps its lower bound.  The finite union of these losses and
the original-law exceptions is negligible compared with
\(M^{-o_n(1)}\).

Measured in the original conditional law given a surviving base
label, the final retained mass of a surviving prefix is at least
\[
p_I M^{-C\epsilon}\ge 2^{-sj_0}M^{-2C\epsilon}.
\]
Dividing the original descendant weight bound by this retained
denominator therefore gives
\[
\Pr(T\in J\mid\text{retained label and prefix})
\le M^{O(\epsilon)}2^{-s(j-j_0)}
\]
for a descendant \(J\) at depth \(j\).  Its support count is bounded
by the preceding Markov truncation.  Parent cells determined by
the label and prefix select entire such fibers.  Finally, making
every truncation against the reference law keeps the absolute
constants from accumulating under successive restrictions; only
a fixed finite union of power-small exceptional sets is removed.
\end{proof}

On a fixed interval of time-depth length \(\tau>0\), the new base is
\(M^{\tau+o_n(1)}\), and the time error in its log units is
\(O(\epsilon_i/\tau)\).  Since
\(\inf_D q(\epsilon,D)\to q_*\), the universal propagation estimate
is therefore available on that interval with exponent
\(q_*-o_i(1)\).  A sufficiently small time-error threshold works
for every finite complexity at an exponent strictly below \(q_*\).
The sufficiently large base threshold may depend on the fixed
complexity, which is all that is needed here.

When the estimate is applied in many parents, fixed-parameter
uniformity supplies the same subpower relative retention and loss
in all of them.  Width complexities and dyadic base ratios that
vary along a sequence, but remain bounded for fixed \(i\), cause
no difficulty: pass to convergent logarithmic ratios, or use a
finite fine power net.  Apply the estimate at a fixed complexity
slightly larger than required.  If this asks for input tests just
below the available floor, round their widths up to that floor;
the arbitrarily small slack changes the input constant by an
arbitrarily small power only.  No application uses a complexity
escaping to infinity along an \(n\)-sequence with \(i\) fixed.

\begin{lemma}[Isolation of terminal widths]\label{geo:fixed-widths}
The failing families may be chosen so that the terminal widths
\(L,w\) are fixed in logarithmic units along each sequence, up to
arbitrarily small prescribed rounding errors.  The orientation
and the two centers remain unrestricted.  The input constant is
still the original \(K\), up to subpower factors, and the time
profile retains error \(O(\epsilon_i)\).
\end{lemma}

\begin{proof}
For each fixed \(i\), choose a finite power net in the permitted
width range, including its endpoints.  Its step tends to zero
as \(i\to\infty\), and is small compared with the margin in the
selected failure.  Round both widths upward.  This enlarges the
test and increases \(W\) by only the allowed small power.

Seek the desired sublaw successively for each pair in this finite
net, always testing against the original \(K\).  If all steps
succeeded, their successive subpower restrictions would give
\eqref{eq:source-3} for every width, contradicting the selected
failure.  Some step therefore fails on a retained law.  Successful
steps need only proceed along a subsequence, and there are finitely
many steps for fixed \(i\).  Lemma~\ref{geo:time-conditioning}
preserves the time profile, and the fixed-power failure margin
absorbs the finite net and normalization losses.
\end{proof}

\subsection{A hairbrush bound and angular uniqueness}

The time profile is now reusable after the stated trims, and the
terminal widths can be isolated.  We next count heavy boxes and
control all their axes through each surviving point.  The proof uses
the intersection-counting mechanism of classical hairbrush arguments;
see Wolff~\cite{Wolff1995Kakeya} and the discussion in
\cite[arXiv version~1, Section~3.7]{TaoVargasVega1998}.
The joint ellipse estimate required here is proved in full below.

Here a \emph{particle} is an outcome of the full joint law, including
any auxiliary labels.  A \emph{lag} is a time depth \(y\in[0,1]\).
For a possible value \(v\) of \(T_y\), independent centers \(u,x\in\R^2\),
and a centered ellipse \(E\), the corresponding box is the event
\[
C_y(v;u,x;E)
=\{T_y=v,\ U\in u+E,\ X_v\in x+M^{-y}E\}.
\]
Its radii are those of \(E\), measured in the variables
\((U,M^yX_v)\), and its mass includes the time-bin probability.
Auxiliary labels do not enter the membership test.  A source box is
the time-free event \(\{U\in u+E,\ V\in x+E\}\).

\Needspace{3\baselineskip}
\begin{lemma}[Hairbrush estimate]\label{geo:hairbrush}
Let \(\mu\) be a subprobability law on the particles.  Fix a lag
and suppose that every box in a chosen time bin, with independently
translated spatial tests of the same ellipse \(F\), satisfies
\(\mu(C)\le H W(F)\).  Fix radii \(L\ge w>0\), and put
\(A=HW(E)\), where \(E\) has radii \(L,w\).  Assume the cap is
available at all radii between \(w\) and \(L\), including bounded
dilations.

Let \(C_0\) be a box of radii \(L,w\).  Suppose each particle
\(z\) in a measurable set \(S\subset C_0\) belongs to an assigned
box \(C_z\) in the same bin, of the same radii, such that
\[
\mu(C_z)\ge A M^{-e},
\qquad
\theta\le\angle(C_z,C_0)\le2\theta,
\qquad
\theta\gtrsim w/L.
\]
Angles are distances between unoriented long axes.  At fixed
polynomial scale complexity, \(m=\mu(S)\) satisfies
\begin{equation}\label{eq:source-5}
m\le M^{2e+o(1)}H W(E)
       \left(\frac{w}{L\theta}\right)^{2\kappa}.
\end{equation}
The same assertion holds for source boxes, with time omitted.
\end{lemma}

\begin{proof}
Replace ellipses by comparable rectangles.  Within the chosen bin,
write \(Z=(Z^{(1)},Z^{(2)})\) for the two spatial variables; at lag
\(y\), these are \((U,M^yX_v)\).  Translations in the two components
remain independent.  Set
\[
\eta=\frac{w}{L\theta}.
\]
The bounded-factor cases with \(\eta\gtrsim1\) follow directly from
\(\mu(S)\le\mu(C_0)\lesssim A\).  We may therefore suppose
\(\eta\lesssim1\).

For \(z,z'\in S\), let \(dL\) be the larger longitudinal separation
of \(Z^{(j)}(z),Z^{(j)}(z')\), measured in the axis of \(C_0\).
For \(d\gtrsim\eta\), the particles \(z'\) at separation at most
\(dL\) lie in a product rectangle of radii \(O(dL),O(w)\).
Consequently their mass is at most
\[
A d^{1+\kappa}.
\]
Suppose \(d\gg\eta\) and \(C_z\cap C_{z'}\ne\varnothing\).
Choose a component realizing this longitudinal separation.  The
transverse separation of its anchors in the axis of \(C_0\) is
\(O(w)\).  Projection onto a normal to the axis of \(C_z\), using
a common point of the two rectangles and their bounded lengths,
gives
\[
dL\theta\lesssim w+L\gamma,
\qquad
\gamma=\angle(C_z,C_{z'}).
\]
Hence \(\gamma\gtrsim d\theta\).

In each component the intersection is contained in a rectangle
with the axis of \(C_z\) and radii
\[
O\!\left(\frac{w}{d\theta}\right)=O(L\eta/d),
\qquad O(w).
\]
The two component intersections use the same orientation, though
their centers may differ.  The assumed cap thus gives
\[
\mu(C_z\cap C_{z'})
\lesssim A(\eta/d)^{1+\kappa}.
\]
Every radius used lies between \(w\) and \(L\), up to bounded
factors.  This argument only uses that a box contains its anchor;
the anchor need not be its center.

For fixed \(z\), the close pairs \(d=O(\eta)\) contribute at most
\(A^2\eta^{1+\kappa}\) to the overlap integral.  Each remaining
dyadic shell contributes at most
\[
\bigl(A d^{1+\kappa}\bigr)
\bigl(A(\eta/d)^{1+\kappa}\bigr)
=A^2\eta^{1+\kappa}.
\]
There are \(O(1+\log(L/w))\) shells, so
\[
\int_S\int_S\mu(C_z\cap C_{z'})\,d\mu(z')\,d\mu(z)
\lesssim mA^2\eta^{1+\kappa}\bigl(1+\log(L/w)\bigr).
\]
The union of the \(C_z\) lies in a product rectangle of radii
\(O(L),O(\theta L)\), of mass at most
\(O(A\eta^{-(1-\kappa)})\).  For
\[
f(p)=\int_S\mathbf1_{C_z}(p)\,d\mu(z),
\]
heaviness and Fubini give
\[
\int f\,d\mu\ge mAM^{-e},
\qquad
\int f^2\,d\mu
=\int_S\int_S\mu(C_z\cap C_{z'})\,d\mu(z')\,d\mu(z).
\]
Cauchy--Schwarz on the enclosing box yields
\[
m^2A^2M^{-2e}
\lesssim A\eta^{-(1-\kappa)}
\,mA^2\eta^{1+\kappa}\bigl(1+\log(L/w)\bigr).
\]
After cancellation this is
\[
m\lesssim A M^{2e}\eta^{2\kappa}\bigl(1+\log(L/w)\bigr).
\]
The logarithm is subpower at fixed polynomial complexity, proving
\eqref{eq:source-5}.  Nothing else in the proof uses time.
\end{proof}

\begin{corollary}[Counting heavy boxes]\label{geo:heavy-count}
Suppose the hypotheses of Lemma~\ref{geo:hairbrush} hold in every
time bin with the same \(H\).  Fix radii \(L,w\), and use an
orientation grid of spacing comparable to \(\rho=w/L\), with
translated grids having bounded overlap at each orientation.
Let \(\mathcal H\) be the family of boxes of mass at least
\(HW(E)M^{-e}\).  Then, over all time bins,
\begin{equation}\label{eq:source-6}
\#\mathcal H
\le \bigl(HW(E)\bigr)^{-1}M^{O_\kappa(e)+o(1)}.
\end{equation}
\end{corollary}

\begin{proof}
Write \(A=HW(E)\) and \(c_\kappa=4/\kappa\).  Choose
\(\tau_M\downarrow0\) sufficiently slowly that all bounded and
logarithmic factors are \(M^{o(\tau_M)}\) and
\(\tau_M\log M\to\infty\).  Set
\[
\Theta_0=\rho M^{c_\kappa(e+\tau_M)}.
\]
From each \(C\in\mathcal H\), remove the particles belonging to
another heavy box whose axis differs from that of \(C\) by more
than \(\Theta_0\).  Summing \eqref{eq:source-5} over angular shells
bounds the removed mass by
\[
A M^{2e-2\kappa c_\kappa(e+\tau_M)+o(\tau_M)}
=A M^{-6e-8\tau_M+o(\tau_M)}
=o(AM^{-e}).
\]
If the cutoff exceeds the angular range there is nothing to
remove.  The remaining set \(G_C\subset C\) has mass at least
\(\tfrac12 AM^{-e}\).

A particle belonging to some \(G_C\) can belong to heavy boxes
only at orientations within \(\Theta_0\) of \(C\).  It therefore
belongs to at most
\[
O(1+\Theta_0/\rho)
\le M^{c_\kappa(e+\tau_M)+o(\tau_M)}
\]
grid boxes: there are this many possible orientations, bounded
translated-grid overlap at each orientation, and only one time
bin.  Summing the masses of the \(G_C\) and using total mass at
most one gives
\[
\tfrac12\,\#\mathcal H\,AM^{-e}
\le M^{c_\kappa(e+\tau_M)+o(\tau_M)}.
\]
In particular,
\[
\#\mathcal H
\le A^{-1}M^{(1+c_\kappa)e+c_\kappa\tau_M+o(\tau_M)},
\]
which proves \eqref{eq:source-6}.
\end{proof}

\Needspace{3\baselineskip}
\begin{corollary}[Pointwise angular exclusion]\label{geo:angular-exclusion}
Let \(\mathcal H\) be the family in
Corollary~\ref{geo:heavy-count}.  Suppose a large sublaw
\(\nu\le\mu\) carries assigned boxes \(C(z)\in\mathcal H\), with
\(z\in C(z)\).  Assume the cap is also available at all enlarged
widths used below.  After removing \(o(\nu(\text{whole space}))\)
mass, the following implication holds pointwise:
\[
z\in M^d D,\quad D\in\mathcal H
\quad\Longrightarrow\quad
\angle(C(z),D)
\le\frac{w}{L}M^{C_\kappa(e+d)+o(1)}.
\]
Here \(D\) is in the particle's time bin, and dilation acts on
each spatial rectangle about its own center.  One retained set
works simultaneously for the dyadic dilations in any fixed
power range \(d\ge0\).
\end{corollary}

\begin{proof}
Keep \(A=HW(E)\), \(\rho=w/L\), \(c_\kappa=4/\kappa\), and take
\[
C_\kappa=\frac{c_\kappa+10}{2\kappa}.
\]
Write \(\nu(\text{whole space})=M^{-r_M}\), with \(r_M\to0\).
Choose \(\tau_M\downarrow0\) slowly enough that
\(r_M=o(\tau_M)\), all counting and logarithmic error exponents
are \(o(\tau_M)\), and \(\tau_M\log M\to\infty\).

Fix a dilation \(Q=M^d\).  Its baseline is
\[
A_Q=HW(QE)=Q^2A=M^{2d}A.
\]
Every enlarged heavy box has mass at least
\[
AM^{-e}=A_QM^{-(e+2d)}.
\]
For each \(C\in\mathcal H\), consider its particles that belong
to some \(QD\), \(D\in\mathcal H\), at angle greater than
\[
\Theta_d=\rho M^{C_\kappa(e+d+\tau_M)}.
\]
Since \(C\subset QC\), apply Lemma~\ref{geo:hairbrush} to the
central box \(QC\), the tilted boxes \(QD\), baseline \(A_Q\),
and tolerance \(e+2d\).  Summing angular shells bounds this
exceptional mass by
\[
A M^{2e+6d-2\kappa C_\kappa(e+d+\tau_M)+o(\tau_M)}.
\]
The application uses the caps between \(Qw\) and \(QL\), with
bounded dilations.

Summing over the original heavy-box list, using the explicit
count in the preceding proof, gives total exceptional mass at
most
\[
\begin{split}
&M^{(3+c_\kappa)e+6d+c_\kappa\tau_M
       -2\kappa C_\kappa(e+d+\tau_M)+o(\tau_M)}\\
&\hspace{2cm}
=M^{-7e-(c_\kappa+4)d-10\tau_M+o(\tau_M)}.
\end{split}
\]
The dyadic dilations in a fixed power range have
\(M^{o(\tau_M)}\) members after the choice of \(\tau_M\).  Their
union has exceptional mass
\[
M^{-10\tau_M+o(\tau_M)}
=o(M^{-r_M}).
\]
Removing this union proves the pointwise implication; its
remaining tolerance \(C_\kappa\tau_M\) is \(o(1)\).
\end{proof}

The heavy list and capacity bounds refer to the fixed prelaw.
The retained set has the stronger pointwise exclusion property for
every relevant high box.  Later conditioning will use this property,
rather than a bound only for a randomly selected pair of boxes.

\subsection{Reduction to bounded spatial complexity}

\begin{lemma}[Bounded complexity]\label{geo:bounded-complexity}
In the contradiction families for Theorem~\ref{geo:main}, the
complexities \(D_i\) may be bounded by a constant depending only
on \(\kappa\).
\end{lemma}

\begin{proof}
Work at the fixed failing widths supplied by
Lemma~\ref{geo:fixed-widths}.  Choose a fixed large integer \(J\),
depending only on \(\kappa\), with \(2\kappa J>30\), and put
\[
f=\min(L,wM^J),\qquad F=(L,f).
\]
Here \(F\) denotes the ellipse of those radii in the orientation
that will be specified for each group.

First suppose \(L>wM^J\).  Call a source box of widths \(L,w\)
\emph{high} if its mass is at least
\[
KW(E)M^{-10}.
\]
Use boundedly enlarged grids so that every terminal test lies
in a source test with these widths up to bounded factors:
from \(U\in E,\ X_v\in M^{-1}E\) we obtain
\(V=X_v-vU\in c+O(E)\).

Consider the particles whose \(P\) lies in two high source boxes
at angular distance greater than \(M^Jw/L\).  Inside any tested
source box, at least one of those two boxes makes angle
\(\gtrsim M^Jw/L\) with the test.  Applying
Lemma~\ref{geo:hairbrush} at the source with \(e=10\), and
summing angular shells, bounds this sublaw's mass in the test by
\[
KW(E)M^{20-2\kappa J+o_n(1)}.
\]
Our choice of \(J\) makes this better than every failing exponent
under consideration.  The sublaw of particles in no high source
box likewise has mass at most \(KW(E)M^{-10}\) in each source
test.  Each of these sublaws therefore satisfies the required
terminal bounds.  Neither can have large mass in the selected
counterexample family.  After their removal, a large sublaw
remains, and all high source boxes through any one surviving
\(P\) have axes within \(O(M^Jw/L)\).

Choose one such axis as a function of \(P\), round it at spacing
comparable to \(f/L=wM^J/L\), and partition by the rounded
direction and by translated \(U,V\) grid boxes of widths \(F\)
in that direction.  This partition depends on \(P\) alone.
When \(L\le wM^J\), instead use isotropic \(F=(L,L)\), with
no angular partition.

Any terminal test that violates the pre-partition threshold
supplies a high source box through all of its retained
particles.  In the eccentric case, their rounded directions
must therefore lie within \(O(f/L)\) of its direction.  There
are boundedly many such rounded directions.  In either case,
the terminal test meets boundedly many translated \(U,V\)
group cells at each allowed direction, because it is contained
in a source \(O(E),O(E)\) test.  Thus each violating terminal
test meets boundedly many groups.  Tests already below the
threshold remain below it under further restrictions.

We now normalize a group.  Let \(A_F\) be the linear map taking
the unit disk to \(F\), and subtract the two group centers
before applying \(A_F^{-1}\).  For a nonsingular linear map
\(A\), extend the notation by \(W(A)=W(AB_2)\).  The singular
value formula is
\[
W(A)=|\det A|^{1-\kappa}\|A\|^{2\kappa},
\]
and consequently
\begin{equation}\label{geo:submultiplicativity}
W(AB)\le W(A)W(B).
\end{equation}
If the group has mass \(m_G\), its normalized input cap is thus
at most
\[
\frac{KW(F)}{m_G}.
\]
This bound holds for every orientation of the normalized test.

For a previously violating terminal test at the selected widths
that meets the group, its target ellipse \(E\) makes angle
\(O(f/L)\) with the group axes, by the preceding partition argument.
If \(A_E\) maps the unit disk to that ellipse, then
\[
\|A_F^{-1}A_E\|\asymp1,
\qquad
|\det(A_F^{-1}A_E)|=w/f.
\]
Its pullback therefore has radii comparable to \(1,w/f\).
Moreover,
\[
W(F)W(A_F^{-1}E)
\lesssim
L^{1+\kappa}f^{1-\kappa}(w/f)^{1-\kappa}
=W(E).
\]
The input floor needed for this target is a constant multiple
of \(M^{-1}w/f\).  Under \(A_F\), the smallest radius of any
such normalized input ellipse is at least
\[
f(M^{-1}w/f)=M^{-1}w\ge M^{-D_i}.
\]
The original input cap is therefore available.  The normalized
complexity is
\[
1+\log_M(f/w)\le J+1.
\]
Bounded coordinate normalizations and endpoint roundings can
be handled with arbitrarily small positive slack in complexity
and the corresponding arbitrarily small loss in the input
constant, as described above.  For all sufficiently large
indices, complexity at most \(J+3\) suffices.

If the limiting worst exponent for complexities bounded by
\(J+3\) were strictly greater than \(q_*\), its universal bound
would improve all these normalized groups beyond the selected
failing threshold.  Fixed-parameter uniformity supplies a
large sublaw in every group with the same subpower losses.
On multiplying back by group masses, the factors \(m_G^{-1}\)
cancel, and the preceding inequality for the capacity weights
returns the original \(KW(E)\).  Recombining uses only boundedly
many groups per previously violating test, while all other
tests already obey the threshold.  The recombined sublaw has
large mass, a contradiction.

The limiting worst exponent at bounded complexities is
therefore \(q_*\), and counterexample families may be chosen
there.  Increasing the fixed bound if necessary proves the
lemma.
\end{proof}

\begin{lemma}[Extension of the width range]\label{geo:range-extension}
After Lemma~\ref{geo:bounded-complexity}, one may extend the
input cap to any prescribed larger bounded power range of
widths, preserve failure at the selected terminal widths, and
place those widths a fixed power below the upper boundary of
the permitted range.
\end{lemma}

\begin{proof}
Put \(\rho=M^{-D_i}\).  Add independent uniform disk noises of
radius \(\rho\) to \(U\) and \(V\), and append those noises to
the base label.  Their independence of the original law
preserves the conditional time profile.

For an ellipse \(E\) of radii \(L,w\), let \(E^+\) have the
same axes and radii \(\max(L,\rho),\max(w,\rho)\).
For one noise variable, the probability of entering a
translated \(E\), given the original point, is at most a
constant times \(|E|/|E^+|\), and vanishes unless that point
lies in a bounded enlargement of the corresponding
translate of \(E^+\).  Independence of the two noises yields
\[
\Pr(U_{\rm sm}\in E,\ V_{\rm sm}\in E)
\lesssim
KW(E^+)\left(\frac{|E|}{|E^+|}\right)^2
\lesssim KW(E).
\]
The final inequality follows because contracting a semiaxis
gets exponent \(2\) from dilution, whereas its exponent in
\(W\) is either \(1+\kappa\) or \(1-\kappa\), both less than
\(2\).  This extends the cap to smaller widths.  Bounded
support extends it to larger widths by clipping the long
radius at a fixed support bound and covering by boundedly
many tests.  If both radii exceed that bound, the total-mass
bound suffices: a fixed covering of the original bounded
support by unit source boxes already gives \(K\gtrsim1\),
with a constant depending only on that support bound.

Failure persists under smoothing.  A good sublaw of the
smoothed joint law pushes back to a sublaw of the original
law with the same mass.  The perturbations of \(U\) and
\(X_v\) are \(O(\rho)\), and the smallest selected terminal
position radius is \(M^{-1}w\ge\rho\).  Thus every original
terminal test maps into boundedly enlarged smoothed tests,
and their good estimates would give the original ones.

Finally, an isotropic power shrinking of both spatial
variables puts the relevant widths a fixed power below the
upper endpoint.  If the shrinking factor is \(B^{-1}\),
the transformed constant is \(KB^2\), while the ellipse
weight becomes \(W(E/B)=B^{-2}W(E)\); their product, and
therefore the failing comparison, is unchanged.  After the
bounded-complexity reduction, all required ranges and
shrinking powers remain bounded.
\end{proof}

From now on we use these extended ranges for universal upper
bounds and angular exclusions.  Factoring through width bands
will use only the original available range, shifted by the
isotropic normalization.  This distinction avoids requiring
an unavailable input cap at a band endpoint.

We have reduced the spatial complexity without changing the
retained-law obstruction.  The backward construction uses two
complementary facts at each selected epoch: a simultaneous cap bounds
all allowed tests, while the finite heavy-box list prevents a large
later sublaw from improving the selected-width bound.

\subsection{Backward construction of heavy ellipses}

Fix a finite dyadic mesh of positive lags \(y\le1\), including
\(1\).  At each lag, apply the universal estimate with exponent
\(q_*-o_i(1)\), propagating from depth \(0\) to depth \(y\).
Lemma~\ref{geo:time-conditioning} and fixed-parameter uniformity
permit these finitely many applications successively.  We
obtain a large normalized \emph{prelaw}, still carrying the
selected failure, with simultaneous bounds
\begin{equation}\label{geo:prelaw-caps}
\Pr(T_y=v,\ U\in E,\ X_v\in M^{-y}E)
\le KW(E)M^{-q_*y+o_i(1)+o_n(1)}
\end{equation}
throughout the required extended width ranges.  Dyadic
roundings are understood up to bounded factors.  For each
fixed mesh, all displayed \(o_i(1)\) errors tend to zero
before the mesh is subsequently refined.

At \(y=1\), call a prelaw box at the selected failing widths
high if its mass is at least
\[
KW(E)M^{-q_*y-e_y},
\]
where \(e_y=o_i(1)>0\) is chosen to dominate the upper-cap
errors and the small margin in the selected failure.
Particles lying in no such box already satisfy a stronger
terminal cap, so they cannot carry large mass.  Assign a
high box to each remaining particle.

Corollary~\ref{geo:heavy-count} bounds the resulting list by
\[
\bigl(KW(E)M^{-q_*y}\bigr)^{-1}
M^{O_\kappa(e_y)+o_n(1)}.
\]
This creates a hereditary obstruction: any later sublaw
carried by the assigned list and satisfying a cap improved
by a sufficiently larger multiple of \(e_y\) has power-small
total mass, simply by summing over the list.  Thus a later
large sublaw cannot make that improvement.  Importantly,
the boxes are high and counted in the prelaw; no later
restriction changes the validity of their covering list.

\Needspace{4\baselineskip}
\begin{lemma}[Backward width band]\label{geo:width-band}
Suppose that at a mesh epoch \(y\) the current large sublaw
has assigned boxes of widths \(L,w\) and the preceding
hereditary obstruction.  Let \(y'<y\) be the previous mesh
epoch.  After a further large restriction, the particles
have assigned high prelaw boxes at \(y'\), with a common
pair of dyadic widths both belonging to
\begin{equation}\label{eq:source-7}
\bigl[wM^{-(y-y')},L\bigr],
\end{equation}
up to bounded endpoint roundings.  The tolerance \(e_{y'}\)
may be chosen to tend to zero sufficiently slowly relative
to \(e_y\) and the propagation errors.
\end{lemma}

\begin{proof}
Call a prelaw box at epoch \(y'\), with both widths in
\eqref{eq:source-7}, high when its mass is at least
\[
KW(F)M^{-q_*y'-e_{y'}}.
\]
Suppose a large sublaw consists of particles belonging to
none of these high boxes.  Every test in this band then
has the improved cap with factor \(M^{-e_{y'}}\): a test
above that threshold would itself be a high prelaw box
containing those particles.  The finite grids and bounded
coverings make this statement simultaneous over all
centers, orientations, and band widths.

Partition this contrary sublaw by the prefix \(v'=T_{y'}\)
and by isotropic parent cells of widths \(L\) for \(U\)
and \(h'L\) for \(X_{v'}\), where
\[
h'=2^{-\lfloor y'n\rfloor}.
\]
Let a parent have centers \(u_Q,x_Q\) and mass \(m_Q\).
Normalize it by
\[
U'=\frac{U-u_Q}{L},\qquad
V'=\frac{X_{v'}-x_Q}{h'L},\qquad
S=\frac{T-v'}{h'}.
\]
The residual time \(S\) lies in \([0,1)\), and its terminal
partition has base
\[
M_{\rm res}
=2^{\lfloor yn\rfloor-\lfloor y'n\rfloor}
=M^{y-y'+o_n(1)}.
\]
If \(s_0\) is the left endpoint of a tested residual interval of
length \(M_{\rm res}^{-1}\), then on that event
\(v=v'+h's_0=T_y\), and
\[
V'+s_0U'
=\frac{X_v-x_Q-(v-v')u_Q}{h'L}.
\]
Both translations are fixed by the parent and this residual prefix.
Moreover, \(h'/M_{\rm res}=2^{-\lfloor yn\rfloor}\asymp M^{-y}\),
so the residual terminal position scale agrees with the original
one, up to the declared dyadic rounding.

A normalized terminal test has ellipse widths \(1,w/L\),
so the smallest input width required for its propagation
is a constant multiple of
\[
M_{\rm res}^{-1}w/L.
\]
In original units this is exactly the lower endpoint
\(wM^{-(y-y')}\) of \eqref{eq:source-7}, up to a bounded
factor.  The band cap hence supplies the full normalized
input bound, with constant
\[
\frac{K L^2M^{-q_*y'-e_{y'}}}{m_Q}
\]
up to subpower factors, because \(W(LF)=L^2W(F)\).
The parent time profile is supplied by
Lemma~\ref{geo:time-conditioning}.  Its complexity is
bounded for the fixed mesh and fixed \(i\).

Apply universal propagation on the residual interval.
Multiplication by the original parent mass cancels
\(m_Q^{-1}\); multiplication by \(W(E/L)\) cancels the
factor \(L^2\).  Thus the resulting sublaw in this parent
satisfies the original terminal test with bound
\[
KW(E)M^{-q_*y-e_{y'}+o_i(1)+o_n(1)}.
\]
The relative retained masses and errors may be taken
uniformly in the parents.

Recombination costs only a bounded factor.  A terminal
bin determines its prefix, and
\[
X_{v'}=X_v+(v'-v)U,\qquad |v'-v|\lesssim M^{-y'}.
\]
As \(U\) varies in a terminal ellipse of long radius \(L\),
the rebased position varies by \(O(M^{-y'}L)\).
The terminal cell consequently meets boundedly many of
the isotropic parent cells in both variables.

Choosing \(e_{y'}\) larger than the relevant multiples of
\(e_y\) and the propagation errors contradicts the
hereditary obstruction at \(y\).  The contrary particles
cannot have large mass.  We may assign high boxes to the
remaining large sublaw and select a common dyadic pair
of widths.  The number of such pairs is subpower, so
this selection retains large mass.  For a fixed finite
mesh the tolerances can be chosen in this backward
order while all still tend to zero.  Bounded rounding
of the band endpoints contributes a vanishing error
in logarithmic units.
\end{proof}

Iterating Lemma~\ref{geo:width-band} backward gives assigned
heavy ellipses at every epoch of the mesh.  The same
list-count obstruction is inherited at each newly assigned
epoch, so the induction can continue.  All restrictions
still compare directly to the prelaw for spatial caps
and to the reference law for the conditional time profile.

\subsection{Width compactness and first-order calibration}

Write the selected radii at epoch \(y\) as
\[
L(y)=M^{-x(y)},\qquad
w(y)=M^{-x(y)-d(y)},\qquad d(y)\ge0,
\]
and write \(\operatorname{axis}_y(z)\) for a particle's
assigned long axis.

The band relation implies, for successive \(y'<y\),
\[
x(y')\ge x(y),\qquad
x(y')+d(y')+y'\le x(y)+d(y)+y,
\]
up to the vanishing rounding errors.  The width depths
remain in a bounded range by
Lemma~\ref{geo:bounded-complexity} and
Lemma~\ref{geo:range-extension}.

Choose successively finer dyadic meshes, with refinement
slow enough that the accumulated errors tend to zero.
On each finite mesh first take sufficiently large \(i\),
then sufficiently large \(n\).  Compactness for the two
monotone functions, or a diagonal argument on dyadic
rationals, gives limiting functions of bounded variation,
still denoted \(x,d\), such that
\[
x\ \text{is nonincreasing},\qquad
x+d+y\ \text{is nondecreasing}.
\]
The finite-stage depths converge at continuity points.  The two
bounded-variation functions have ordinary derivatives almost
everywhere; we use the monotone and bounded-variation differentiation
theorems in \cite[Theorem~1.6.25 and Corollary~1.6.35]{Tao2011Measure}.
This differentiability concerns the limiting path.  The finite paths
will approximate its first-order window by the separate choices in
Lemma~\ref{geo:calibration}.

\Needspace{4\baselineskip}
\begin{lemma}[Stability of neighboring axes]\label{geo:axis-stability}
The large sublaws may be chosen so that neighboring mesh
epochs satisfy
\begin{equation}\label{eq:source-8}
\begin{split}
\angle\bigl(\operatorname{axis}_{y}(z),
            \operatorname{axis}_{y'}(z)\bigr)
&\le M^{-d(y)+C_\kappa\Delta+o(1)},\\
\Delta&=|y'-y|+|x(y')-x(y)|\\
&\hspace{1.5cm}
+|(x+d)(y')-(x+d)(y)|.
\end{split}
\end{equation}
They also satisfy the pointwise single-epoch exclusions of
Corollary~\ref{geo:angular-exclusion}, including the allowed
dilations.  These assertions are needed only when
\(\Delta\) is small.
\end{lemma}

\begin{proof}
Let \(\mu\) be the prelaw and \(y'<y\).  Put
\[
\begin{split}
h&=y-y',\qquad a_0=|x(y)-x(y')|,\\
b_0&=|(x+d)(y)-(x+d)(y')|,\\
B_0&=(1+\kappa)a_0+(1-\kappa)b_0.
\end{split}
\]
A later box has a single earlier prefix \(v'\), and
\[
M^{y'}X_{v'}=M^{-h}M^yX_v+M^{y'}(v'-v)U.
\]
The last coefficient is bounded, so its particles fit an earlier-bin
box with the same axis and the later widths, up to bounded factors.

Let \(F\) have radii \(L_F=\max(L(y),L(y'))\) and
\(w_F=\max(w(y),w(y'))\).  Enlarge the transported box and the
earlier assigned box along their own axes to these common widths.
For \(j\in\{y,y'\}\), let \(E_j\) have radii \(L(j),w(j)\), and set
\[
A_j=KW(E_j)M^{-q_*j},\qquad
A_F=KW(F)M^{-q_*y'}.
\]
The capacity ratios and the later box's original heaviness give,
up to the existing vanishing errors,
\[
\frac{A_F}{A_{y'}}\le M^{B_0},\qquad
\mu(\text{transported later box})
 \ge A_FM^{-e_y-q_*h-B_0},\qquad
\frac{w_F}{L_F}\le M^{-d(y)+b_0}.
\]
The extended width range supplies the caps at these common widths
and at the intermediate widths required by Lemma~\ref{geo:hairbrush}.

Apply that lemma inside each enlarged earlier box, assigning to
each exceptional particle its transported later box.  Use angular
cutoff \(M^{-d(y)+C_\kappa\Delta+\sigma}\), where \(\sigma>0\)
will tend to zero, and sum the angular shells.  Sum next over the
\emph{original} earlier assigned list, whose count is at most
\(A_{y'}^{-1}M^{(1+4/\kappa)e_{y'}+o(1)}\) by the proof of
Corollary~\ref{geo:heavy-count}.  The factor \(A_F\) in the
hairbrush bound cancels this list baseline up to
\(A_F/A_{y'}\le M^{B_0}\).  Thus the total exceptional mass is at
most
\[
M^{3B_0+2q_*h+2\kappa b_0+(1+4/\kappa)e_{y'}+2e_y
      -2\kappa C_\kappa\Delta-2\kappa\sigma+o(1)}.
\]
Since \(B_0,h,b_0=O(\Delta)\), a sufficiently large
\(C_\kappa\) absorbs their contributions.  Write the current
retained mass as \(M^{-r_M}\), with \(r_M=o_n(1)\).  At each
fixed mesh, choose \(\sigma=\sigma_i\downarrow0\) slowly enough
that \(e_y,e_{y'}\) and all \(o_i(1)\) errors are \(o(\sigma_i)\).
With \(i\) fixed, take \(n\) sufficiently large that \(r_M\), the
\(o_n(1)\) errors, and logarithmic counting losses are negligible
compared with \(\sigma_i\), while \(\sigma_i\log M\to\infty\).
The exceptional mass is then \(o(M^{-r_M})\), as required.

Perform these removals for each neighboring pair and apply
Corollary~\ref{geo:angular-exclusion} at each epoch, including its
dyadic dilation grid.  At each fixed mesh the number of requirements
is subpower.  Choose the vanishing margins also to dominate the
logarithm of this number divided by \(\log M\); the union of the
exceptional sets still has negligible relative mass.
\end{proof}

\Needspace{4\baselineskip}
\begin{lemma}[Calibrated widths and a reference axis]
\label{geo:calibration}
Let \(y_0\in(0,1)\) be a point of ordinary differentiability
of both limiting width depths, and put
\[
x_0=x(y_0),\qquad d_0=d(y_0),\qquad
a=-x'(y_0),\qquad b=-(x+d)'(y_0),\qquad
\lambda=a-b.
\]
Then \(a\ge0\), \(b\le1\), and
\(\lambda\ne0\) implies \(d_0>0\).

There are thicknesses \(\delta\downarrow0\), stretches
\(H_\delta\uparrow\infty\), and finite-stage families
whose local mesh epochs \(y=y_0+\delta v\),
\(|v|\le H_\delta\), satisfy uniformly
\begin{equation}\label{geo:affine-widths}
\begin{split}
x(y)&=x_0-a\delta v+o(\delta),\\
x(y)+d(y)&=x_0+d_0-b\delta v+o(\delta).
\end{split}
\end{equation}
Their neighboring increments in \eqref{eq:source-8}
have \(\Delta=o(\delta)\).  All cap and list errors are
\(o(\delta)\), the mesh spacing is \(o(\delta)\), and
\(\log n/n=o(\delta)\), while \(\delta\log M\to\infty\).

For each particle, choose as reference the assigned
axis at the end of the stretched interval having
the deeper affine aspect.  Then
\begin{equation}\label{geo:reference-axis}
\angle\bigl(\operatorname{axis}_{\rm ref}(z),
            \operatorname{axis}_{y_0+\delta v}(z)\bigr)
\le M^{-d_0-\delta\lambda v+o(\delta)}
\end{equation}
uniformly on the local mesh.  If \(\lambda=0\), any
one local epoch may be used as reference.  When
\(d_0>0\), a further constant-mass restriction permits
the reference directions to be represented by
\((1,R^\circ)\), with \(R^\circ\) uniformly bounded.
\end{lemma}

\begin{proof}
The two monotonicities give \(a\ge0\) and \(b\le1\).
Also \(d'=\lambda\).  Since \(d\ge0\), a differentiable
interior point with \(d_0=0\) is a local minimum and
has derivative zero.  This proves the assertion
when \(\lambda\ne0\).

Differentiability gives a common modulus
\(\omega(h)\to0\) for the two depth functions, with
affine approximation error at most
\(|h|\omega(|h|)\).  Choose \(H_\delta\to\infty\)
so slowly that \(\delta H_\delta\to0\) and
\[
H_\delta\omega(\delta H_\delta)\to0.
\]
The errors are then \(o(\delta)\) uniformly for
\(|h|\le\delta H_\delta\).

To transfer this conclusion to the finite-stage
families, first fix \(\delta\) and a finite accuracy
grid of continuity points.  Convergence at these
points, together with the monotonicity of \(x\)
and \(x+d+y\), traps intermediate width values
between nearby ones.  Next take sufficiently
advanced meshes, sufficiently large \(i\), and
sufficiently large \(n\).  A diagonal choice gives
\eqref{geo:affine-widths}, makes every cap, list,
and rounding error \(o(\delta)\), and makes the
neighboring depth increments \(o(\delta)\).  For the angular
bounds, first fix \(\delta\) and prescribe a margin \(o(\delta)\).
Choose the mesh and \(i\) so that \(\sigma_i\) lies below that
margin and every \(e_y\) and outer error is \(o(\sigma_i)\); then
choose \(n\) so that \(r_M\), the finite-base errors, and
\(\log n/n\) are \(o(\sigma_i)\).  Choose the single-epoch
angular-exclusion margins with the same hierarchy.
We may also impose
\(\log n/n=o(\delta)\),
\(\delta\log M\to\infty\), and a number of epochs
subpower in \(M^\delta\).

If \(\lambda>0\), use the latest local epoch as
reference; if \(\lambda<0\), use the earliest.
Along the chain from a given epoch to that
reference, every intermediate aspect depth is
at least \(d_0+\delta\lambda v-o(\delta)\),
by \eqref{geo:affine-widths}.  Summing
\eqref{eq:source-8} along this chain gives
\eqref{geo:reference-axis}: the worst angular
error determines the logarithmic scale, and the
number of links costs \(M^{o(\delta)}\).
For \(\lambda=0\), all aspect depths are
\(d_0+o(\delta)\), so the same argument works
with any reference epoch.

Finally cover the space of unoriented directions
by finitely many sectors.  On one sector a large
constant fraction of the law remains; after a
fixed rotation, its slopes are uniformly bounded.
This gives the stated representation of the
reference axis.
\end{proof}

\Needspace{4\baselineskip}
\begin{corollary}[Pointwise axis prediction]\label{geo:axis-prediction}
In the families of Lemma~\ref{geo:calibration}, two
retained particles with the same \(P=(U,V)\) and the
same time prefix at a local epoch have reference
axes agreeing to that epoch's aspect precision,
up to a factor \(M^{o(\delta)}\).
In a bounded-slope sector, the number of possible
reference-slope bins at depth
\(d_0+\delta\lambda v\), given the base label and
the time prefix at \(y_0+\delta v\), is at most
\(M^{o(\delta)}\).
\end{corollary}

\begin{proof}
The two particles have exactly the same displayed
point \((U,M^yX_{T_y})\) at the epoch in question.
Each therefore lies in the other's assigned high
prelaw box.  The pointwise exclusion of
Corollary~\ref{geo:angular-exclusion} bounds the
distance between their assigned axes by the
aspect ratio times \(M^{o(\delta)}\).
Equation~\eqref{geo:reference-axis} transfers the
same bound to their reference axes.  Angle and
slope are comparable in the fixed sector.

This is a pointwise membership argument against
all high prelaw boxes.  It does not require the
two particles to carry the same auxiliary label,
or their coincidence to have positive pair
probability.  Since the base label determines
\(P\), the asserted conditional bin count follows.
\end{proof}

The geometric reduction supplies a first-order width path and a
reference axis with subpower error in the local logarithmic unit
\(\delta\log M\).  The exact line and time prefix predict that axis
at the aspect precision of each local epoch.  The capacity bounds,
heavy-box lists, and pointwise angular exclusions remain available
after any further restriction of relative mass
\(\exp(-o(\delta\log M))\); normalization then costs
\(\exp(o(\delta\log M))\).  For each fixed finite list of later
restrictions whose relative masses are subpower with \(\delta\)
fixed, choose the sufficiently large finite stage before decreasing
\(\delta\), so their normalization losses are negligible in this
local unit.
The next section records these properties in finite measured states
before passing to entropy limits.
